% Options for packages loaded elsewhere
% Options for packages loaded elsewhere
\PassOptionsToPackage{unicode}{hyperref}
\PassOptionsToPackage{hyphens}{url}
\PassOptionsToPackage{dvipsnames,svgnames,x11names}{xcolor}
%
\documentclass[
  ngerman,
  letterpaper,
  DIV=11]{scrreprt}
\usepackage{xcolor}
\usepackage{amsmath,amssymb}
\setcounter{secnumdepth}{-\maxdimen} % remove section numbering
\usepackage{iftex}
\ifPDFTeX
  \usepackage[T1]{fontenc}
  \usepackage[utf8]{inputenc}
  \usepackage{textcomp} % provide euro and other symbols
\else % if luatex or xetex
  \usepackage{unicode-math} % this also loads fontspec
  \defaultfontfeatures{Scale=MatchLowercase}
  \defaultfontfeatures[\rmfamily]{Ligatures=TeX,Scale=1}
\fi
\usepackage{lmodern}
\ifPDFTeX\else
  % xetex/luatex font selection
\fi
% Use upquote if available, for straight quotes in verbatim environments
\IfFileExists{upquote.sty}{\usepackage{upquote}}{}
\IfFileExists{microtype.sty}{% use microtype if available
  \usepackage[]{microtype}
  \UseMicrotypeSet[protrusion]{basicmath} % disable protrusion for tt fonts
}{}
\makeatletter
\@ifundefined{KOMAClassName}{% if non-KOMA class
  \IfFileExists{parskip.sty}{%
    \usepackage{parskip}
  }{% else
    \setlength{\parindent}{0pt}
    \setlength{\parskip}{6pt plus 2pt minus 1pt}}
}{% if KOMA class
  \KOMAoptions{parskip=half}}
\makeatother
% Make \paragraph and \subparagraph free-standing
\makeatletter
\ifx\paragraph\undefined\else
  \let\oldparagraph\paragraph
  \renewcommand{\paragraph}{
    \@ifstar
      \xxxParagraphStar
      \xxxParagraphNoStar
  }
  \newcommand{\xxxParagraphStar}[1]{\oldparagraph*{#1}\mbox{}}
  \newcommand{\xxxParagraphNoStar}[1]{\oldparagraph{#1}\mbox{}}
\fi
\ifx\subparagraph\undefined\else
  \let\oldsubparagraph\subparagraph
  \renewcommand{\subparagraph}{
    \@ifstar
      \xxxSubParagraphStar
      \xxxSubParagraphNoStar
  }
  \newcommand{\xxxSubParagraphStar}[1]{\oldsubparagraph*{#1}\mbox{}}
  \newcommand{\xxxSubParagraphNoStar}[1]{\oldsubparagraph{#1}\mbox{}}
\fi
\makeatother

\usepackage{color}
\usepackage{fancyvrb}
\newcommand{\VerbBar}{|}
\newcommand{\VERB}{\Verb[commandchars=\\\{\}]}
\DefineVerbatimEnvironment{Highlighting}{Verbatim}{commandchars=\\\{\}}
% Add ',fontsize=\small' for more characters per line
\usepackage{framed}
\definecolor{shadecolor}{RGB}{241,243,245}
\newenvironment{Shaded}{\begin{snugshade}}{\end{snugshade}}
\newcommand{\AlertTok}[1]{\textcolor[rgb]{0.68,0.00,0.00}{#1}}
\newcommand{\AnnotationTok}[1]{\textcolor[rgb]{0.37,0.37,0.37}{#1}}
\newcommand{\AttributeTok}[1]{\textcolor[rgb]{0.40,0.45,0.13}{#1}}
\newcommand{\BaseNTok}[1]{\textcolor[rgb]{0.68,0.00,0.00}{#1}}
\newcommand{\BuiltInTok}[1]{\textcolor[rgb]{0.00,0.23,0.31}{#1}}
\newcommand{\CharTok}[1]{\textcolor[rgb]{0.13,0.47,0.30}{#1}}
\newcommand{\CommentTok}[1]{\textcolor[rgb]{0.37,0.37,0.37}{#1}}
\newcommand{\CommentVarTok}[1]{\textcolor[rgb]{0.37,0.37,0.37}{\textit{#1}}}
\newcommand{\ConstantTok}[1]{\textcolor[rgb]{0.56,0.35,0.01}{#1}}
\newcommand{\ControlFlowTok}[1]{\textcolor[rgb]{0.00,0.23,0.31}{\textbf{#1}}}
\newcommand{\DataTypeTok}[1]{\textcolor[rgb]{0.68,0.00,0.00}{#1}}
\newcommand{\DecValTok}[1]{\textcolor[rgb]{0.68,0.00,0.00}{#1}}
\newcommand{\DocumentationTok}[1]{\textcolor[rgb]{0.37,0.37,0.37}{\textit{#1}}}
\newcommand{\ErrorTok}[1]{\textcolor[rgb]{0.68,0.00,0.00}{#1}}
\newcommand{\ExtensionTok}[1]{\textcolor[rgb]{0.00,0.23,0.31}{#1}}
\newcommand{\FloatTok}[1]{\textcolor[rgb]{0.68,0.00,0.00}{#1}}
\newcommand{\FunctionTok}[1]{\textcolor[rgb]{0.28,0.35,0.67}{#1}}
\newcommand{\ImportTok}[1]{\textcolor[rgb]{0.00,0.46,0.62}{#1}}
\newcommand{\InformationTok}[1]{\textcolor[rgb]{0.37,0.37,0.37}{#1}}
\newcommand{\KeywordTok}[1]{\textcolor[rgb]{0.00,0.23,0.31}{\textbf{#1}}}
\newcommand{\NormalTok}[1]{\textcolor[rgb]{0.00,0.23,0.31}{#1}}
\newcommand{\OperatorTok}[1]{\textcolor[rgb]{0.37,0.37,0.37}{#1}}
\newcommand{\OtherTok}[1]{\textcolor[rgb]{0.00,0.23,0.31}{#1}}
\newcommand{\PreprocessorTok}[1]{\textcolor[rgb]{0.68,0.00,0.00}{#1}}
\newcommand{\RegionMarkerTok}[1]{\textcolor[rgb]{0.00,0.23,0.31}{#1}}
\newcommand{\SpecialCharTok}[1]{\textcolor[rgb]{0.37,0.37,0.37}{#1}}
\newcommand{\SpecialStringTok}[1]{\textcolor[rgb]{0.13,0.47,0.30}{#1}}
\newcommand{\StringTok}[1]{\textcolor[rgb]{0.13,0.47,0.30}{#1}}
\newcommand{\VariableTok}[1]{\textcolor[rgb]{0.07,0.07,0.07}{#1}}
\newcommand{\VerbatimStringTok}[1]{\textcolor[rgb]{0.13,0.47,0.30}{#1}}
\newcommand{\WarningTok}[1]{\textcolor[rgb]{0.37,0.37,0.37}{\textit{#1}}}

\usepackage{longtable,booktabs,array}
\usepackage{calc} % for calculating minipage widths
% Correct order of tables after \paragraph or \subparagraph
\usepackage{etoolbox}
\makeatletter
\patchcmd\longtable{\par}{\if@noskipsec\mbox{}\fi\par}{}{}
\makeatother
% Allow footnotes in longtable head/foot
\IfFileExists{footnotehyper.sty}{\usepackage{footnotehyper}}{\usepackage{footnote}}
\makesavenoteenv{longtable}
\usepackage{graphicx}
\makeatletter
\newsavebox\pandoc@box
\newcommand*\pandocbounded[1]{% scales image to fit in text height/width
  \sbox\pandoc@box{#1}%
  \Gscale@div\@tempa{\textheight}{\dimexpr\ht\pandoc@box+\dp\pandoc@box\relax}%
  \Gscale@div\@tempb{\linewidth}{\wd\pandoc@box}%
  \ifdim\@tempb\p@<\@tempa\p@\let\@tempa\@tempb\fi% select the smaller of both
  \ifdim\@tempa\p@<\p@\scalebox{\@tempa}{\usebox\pandoc@box}%
  \else\usebox{\pandoc@box}%
  \fi%
}
% Set default figure placement to htbp
\def\fps@figure{htbp}
\makeatother



\ifLuaTeX
\usepackage[bidi=basic]{babel}
\else
\usepackage[bidi=default]{babel}
\fi
% get rid of language-specific shorthands (see #6817):
\let\LanguageShortHands\languageshorthands
\def\languageshorthands#1{}
\ifLuaTeX
  \usepackage[german]{selnolig} % disable illegal ligatures
\fi


\setlength{\emergencystretch}{3em} % prevent overfull lines

\providecommand{\tightlist}{%
  \setlength{\itemsep}{0pt}\setlength{\parskip}{0pt}}



 


\usepackage{my_quarto_tools}
\usepackage{enumitem}
\usepackage{changepage}
\setlist[description]{style=nextline, leftmargin=2.5em}
\newcommand{\cmd}[1]{\texttt{#1}\xspace}
\KOMAoption{captions}{tableheading}
\makeatletter
\@ifpackageloaded{caption}{}{\usepackage{caption}}
\AtBeginDocument{%
\ifdefined\contentsname
  \renewcommand*\contentsname{Inhaltsverzeichnis}
\else
  \newcommand\contentsname{Inhaltsverzeichnis}
\fi
\ifdefined\listfigurename
  \renewcommand*\listfigurename{Abbildungsverzeichnis}
\else
  \newcommand\listfigurename{Abbildungsverzeichnis}
\fi
\ifdefined\listtablename
  \renewcommand*\listtablename{Tabellenverzeichnis}
\else
  \newcommand\listtablename{Tabellenverzeichnis}
\fi
\ifdefined\figurename
  \renewcommand*\figurename{Abbildung}
\else
  \newcommand\figurename{Abbildung}
\fi
\ifdefined\tablename
  \renewcommand*\tablename{Tabelle}
\else
  \newcommand\tablename{Tabelle}
\fi
}
\@ifpackageloaded{float}{}{\usepackage{float}}
\floatstyle{ruled}
\@ifundefined{c@chapter}{\newfloat{codelisting}{h}{lop}}{\newfloat{codelisting}{h}{lop}[chapter]}
\floatname{codelisting}{Listing}
\newcommand*\listoflistings{\listof{codelisting}{Listingverzeichnis}}
\makeatother
\makeatletter
\makeatother
\makeatletter
\@ifpackageloaded{caption}{}{\usepackage{caption}}
\@ifpackageloaded{subcaption}{}{\usepackage{subcaption}}
\makeatother
\usepackage{bookmark}
\IfFileExists{xurl.sty}{\usepackage{xurl}}{} % add URL line breaks if available
\urlstyle{same}
\hypersetup{
  pdftitle={Maria-DB},
  pdfauthor={Wolfgang Höfer},
  pdflang={de},
  colorlinks=true,
  linkcolor={blue},
  filecolor={Maroon},
  citecolor={Blue},
  urlcolor={Blue},
  pdfcreator={LaTeX via pandoc}}


\title{Maria-DB}
\author{Wolfgang Höfer}
\date{}
\begin{document}
\maketitle


\newcounter{anr}
\setcounter{anr}{1}

\section{Reines Mariadb}\label{reines-mariadb}

Hier wird die Erstellung mittels Dockerfile beschrieben. Im Kurs-Handout
kommt NUR \texttt{docker\ compose} zum Einsatz!

\subsection{Dockerfile erstellen}\label{dockerfile-erstellen}

Ordner und Dockerfile erstellen

\begin{Shaded}
\begin{Highlighting}[]
\FunctionTok{mkdir}\NormalTok{ maria}
\BuiltInTok{cd}\NormalTok{ maria}
\FunctionTok{nano}\NormalTok{ Dockerfile}
\end{Highlighting}
\end{Shaded}

Alle Einrückungen in dieser Datei müssen mit Leerzeichen und nicht mit
Tabulatoren erfolgen! Zum einfacheren Zählen deute ich Leerzeichen
\textbf{am Anfang der Zeile} mit Punkten an.

\textbf{Image festlegen}

\begin{Shaded}
\begin{Highlighting}[]
\ExtensionTok{FROM}\NormalTok{ mariadb:11   }\CommentTok{\# vermeide "latest"}
\end{Highlighting}
\end{Shaded}

\textbf{Variablen festlegen}\\
(Können auch manuell in Umgebungsvariablen geschrieben werden, wenn
Sicherheit ein Thema ist. )

\begin{Shaded}
\begin{Highlighting}[]
\CommentTok{\# ENV Variablen für MariaDB}
\ExtensionTok{ENV}\NormalTok{ MARIADB\_ROOT\_PASSWORD=rootpass}
\ExtensionTok{ENV}\NormalTok{ MARIADB\_DATABASE=mydatabase}
\ExtensionTok{ENV}\NormalTok{ MARIADB\_USER=myuser}
\ExtensionTok{ENV}\NormalTok{ MARIADB\_PASSWORD=mypassword}
\end{Highlighting}
\end{Shaded}

\textbf{Speicherort der Daten}\\
Info für Docker, welcher Ordner externe Daten enthält

\begin{Shaded}
\begin{Highlighting}[]
\CommentTok{\# Datenverzeichnis}
\ExtensionTok{VOLUME} \PreprocessorTok{[}\StringTok{"/var/lib/mysql"}\PreprocessorTok{]}
\end{Highlighting}
\end{Shaded}

\textbf{Nur zur Doku}

\begin{Shaded}
\begin{Highlighting}[]
\CommentTok{\# Standard MariaDB Port}
\ExtensionTok{EXPOSE}\NormalTok{ 3306}
\end{Highlighting}
\end{Shaded}

\subsection{Erstellen des Images}\label{erstellen-des-images}

\begin{Shaded}
\begin{Highlighting}[]
\ExtensionTok{docker}\NormalTok{ build }\AttributeTok{{-}t}\NormalTok{ my{-}mariadb .   }\CommentTok{\# Punkt ist wichtig!}
\end{Highlighting}
\end{Shaded}

\subsection{Starten des Containers}\label{starten-des-containers}

\begin{Shaded}
\begin{Highlighting}[]
\ExtensionTok{docker}\NormalTok{ run }\AttributeTok{{-}d} \DataTypeTok{\textbackslash{}}
  \AttributeTok{{-}{-}name}\NormalTok{ mariadb{-}container }\DataTypeTok{\textbackslash{}}
  \AttributeTok{{-}p}\NormalTok{ 3333:3306 }\DataTypeTok{\textbackslash{}}
  \AttributeTok{{-}v} \VariableTok{$(}\BuiltInTok{pwd}\VariableTok{)}\NormalTok{/data:/var/lib/mysql }\DataTypeTok{\textbackslash{}}
\NormalTok{  my{-}mariadb}
\end{Highlighting}
\end{Shaded}

\subsection{Erweiterung auf externes
Config-File}\label{erweiterung-auf-externes-config-file}

Ergänzung Dockerfile

\begin{Shaded}
\begin{Highlighting}[]
\ExtensionTok{VOLUME} \PreprocessorTok{[}\StringTok{"/etc/mysql/conf.d"}\PreprocessorTok{]}
\end{Highlighting}
\end{Shaded}

Image neu bauen

\begin{Shaded}
\begin{Highlighting}[]
\ExtensionTok{docker}\NormalTok{ build }\AttributeTok{{-}t}\NormalTok{ my{-}mariadb .   }\CommentTok{\# Punkt ist wichtig!}
\end{Highlighting}
\end{Shaded}

Neuer Startaufruf

\begin{Shaded}
\begin{Highlighting}[]
\ExtensionTok{docker}\NormalTok{ run }\AttributeTok{{-}d} \DataTypeTok{\textbackslash{}}
  \AttributeTok{{-}{-}name}\NormalTok{ mariadb{-}container }\DataTypeTok{\textbackslash{}}
  \AttributeTok{{-}p}\NormalTok{ 3333:3306 }\DataTypeTok{\textbackslash{}}
  \AttributeTok{{-}v} \VariableTok{$(}\BuiltInTok{pwd}\VariableTok{)}\NormalTok{/data:/var/lib/mysql }\DataTypeTok{\textbackslash{}}
  \AttributeTok{{-}v} \VariableTok{$(}\BuiltInTok{pwd}\VariableTok{)}\NormalTok{/my.cnf:/etc/mysql/conf.d/my.cnf }\DataTypeTok{\textbackslash{}}
\NormalTok{  my{-}mariadb}
\end{Highlighting}
\end{Shaded}

\section{MariaDB mit phpmyadmin}\label{mariadb-mit-phpmyadmin}

\textbf{Prinzip:} Ein Dienst pro Container

Verwendung von Docker-Compose zur Verwaltung mehrerer Container.

\subsection{}\label{section}

mkdir mariadb-stack cd mariadb-stack

mkdir data

mkdir config

\chapter{Compose-File}\label{compose-file}

(Jeder Level der Einrückung sind 2 Leerzeichen) \# dockercompose.yaml

\begin{Shaded}
\begin{Highlighting}[]
\ExtensionTok{version:} \StringTok{"3.9"}

\ExtensionTok{services:}
  \ExtensionTok{mariadb:}
    \ExtensionTok{image:}\NormalTok{ mariadb:11}
    \ExtensionTok{container\_name:}\NormalTok{ mariadb}
    \ExtensionTok{restart:}\NormalTok{ unless{-}stopped}

    \ExtensionTok{ports:}
      \ExtensionTok{{-}} \StringTok{"3333:3306"}

    \ExtensionTok{environment:}
      \ExtensionTok{MARIADB\_ROOT\_PASSWORD:}\NormalTok{ rootpass}
      \ExtensionTok{MARIADB\_DATABASE:}\NormalTok{ mydatabase}
      \ExtensionTok{MARIADB\_USER:}\NormalTok{ myuser}
      \ExtensionTok{MARIADB\_PASSWORD:}\NormalTok{ mypassword}

    \ExtensionTok{volumes:}
      \ExtensionTok{{-}}\NormalTok{ ./data:/var/lib/mysql}
      
      
  \ExtensionTok{phpmyadmin:}
    \ExtensionTok{image:}\NormalTok{ phpmyadmin:latest}
    \ExtensionTok{container\_name:}\NormalTok{ phpmyadmin}
    \ExtensionTok{restart:}\NormalTok{ unless{-}stopped}
    \ExtensionTok{ports:}
      \ExtensionTok{{-}} \StringTok{"8080:80"}
    \ExtensionTok{environment:}
      \ExtensionTok{PMA\_HOST:}\NormalTok{ mariadb}
      \ExtensionTok{PMA\_PORT:}\NormalTok{ 3306}

    \ExtensionTok{volumes:}
      \ExtensionTok{{-}}\NormalTok{ ./data:/var/lib/mysql}
      \ExtensionTok{{-}}\NormalTok{ ./config/my.cnf:/etc/mysql/conf.d/my.cnf}
\end{Highlighting}
\end{Shaded}

nano config/my.cnf

{[}mysqld{]} sql\_mode=STRICT\_TRANS\_TABLES

\chapter{Mehrere Instanzen}\label{mehrere-instanzen}

Bei groß dimensionierter Hardware kann man auch überlegen,
Arbeitsgruppen einen eigenen Server zu \emph{spendieren}. Hier ist es
aber wirklich wichtig, viel RAM und CPUs zu besitzen. Die Gruppen werden
nachfolgend als USER bezeichnet.

\textbf{compose-template.yml}

\begin{Shaded}
\begin{Highlighting}[]
\ExtensionTok{version:} \StringTok{"3.9"}

\ExtensionTok{services:}
  \ExtensionTok{mariadb:}
    \ExtensionTok{image:}\NormalTok{ mariadb:11}
    \ExtensionTok{container\_name:}\NormalTok{ mariadb\_}\VariableTok{$\{USER\}}
    \ExtensionTok{restart:}\NormalTok{ unless{-}stopped}
    \ExtensionTok{ports:}
      \ExtensionTok{{-}} \StringTok{"}\VariableTok{$\{PORT\}}\StringTok{:3306"}
    \ExtensionTok{environment:}
      \ExtensionTok{MARIADB\_ROOT\_PASSWORD:} \VariableTok{$\{ROOTPASS\}}
      \ExtensionTok{MARIADB\_DATABASE:} \VariableTok{$\{DB\}}
      \ExtensionTok{MARIADB\_USER:} \VariableTok{$\{USER\}}
      \ExtensionTok{MARIADB\_PASSWORD:} \VariableTok{$\{PASS\}}
    \ExtensionTok{volumes:}
      \ExtensionTok{{-}}\NormalTok{ ./data:/var/lib/mysql}
      \ExtensionTok{{-}}\NormalTok{ ./config/my.cnf:/etc/mysql/conf.d/my.cnf}

  \ExtensionTok{phpmyadmin:}
    \ExtensionTok{image:}\NormalTok{ phpmyadmin:latest}
    \ExtensionTok{container\_name:}\NormalTok{ phpmyadmin\_}\VariableTok{$\{USER\}}
    \ExtensionTok{restart:}\NormalTok{ unless{-}stopped}
    \ExtensionTok{ports:}
      \ExtensionTok{{-}} \StringTok{"}\VariableTok{$\{PMA\_PORT\}}\StringTok{:80"}
    \ExtensionTok{environment:}
      \ExtensionTok{PMA\_HOST:}\NormalTok{ mariadb\_}\VariableTok{$\{USER\}}
      \ExtensionTok{PMA\_PORT:}\NormalTok{ 3306}
\end{Highlighting}
\end{Shaded}

Nachfolgendes Script \texttt{create\_server.sh} kann dann für jede
Gruppe aufgerufen werden:

\begin{Shaded}
\begin{Highlighting}[]
\CommentTok{\#!/bin/bash}

\VariableTok{USER}\OperatorTok{=}\VariableTok{$1}
\VariableTok{PORT}\OperatorTok{=}\VariableTok{$2}
\VariableTok{PMA\_PORT}\OperatorTok{=}\VariableTok{$3}

\FunctionTok{mkdir} \AttributeTok{{-}p} \VariableTok{$USER}\NormalTok{/}\DataTypeTok{\{data}\OperatorTok{,}\DataTypeTok{config\}}

\FunctionTok{cp}\NormalTok{ compose{-}template.yml }\VariableTok{$USER}\NormalTok{/docker{-}compose.yml}

\FunctionTok{cat} \OperatorTok{\textgreater{}} \VariableTok{$USER}\NormalTok{/.env }\OperatorTok{\textless{}\textless{}EOF}
\StringTok{USER=}\VariableTok{$USER}
\StringTok{PORT=}\VariableTok{$PORT}
\StringTok{PMA\_PORT=}\VariableTok{$PMA\_PORT}
\StringTok{ROOTPASS=rootpass\_}\VariableTok{$\{USER\}}
\StringTok{DB=db\_}\VariableTok{$\{USER\}}
\StringTok{PASS=pass\_}\VariableTok{$\{USER\}}
\OperatorTok{EOF}

\BuiltInTok{cd} \VariableTok{$USER}
\ExtensionTok{docker}\NormalTok{ compose up }\AttributeTok{{-}d}
\end{Highlighting}
\end{Shaded}

\textbf{Verwendung}

\begin{Shaded}
\begin{Highlighting}[]
\FunctionTok{bash}\NormalTok{ create{-}server.sh gruppe\_1 3333 8080}
\FunctionTok{bash}\NormalTok{ create{-}server.sh gruppe\_2 3334 8081}
\end{Highlighting}
\end{Shaded}





\end{document}
